\section{Hands-On Practical Lab: Building Embedded Linux with Buildroot}

This lab demonstrates how to construct, cross-compile, flash, and boot a minimalist, custom embedded Linux distribution targeting the \textbf{STMicroelectronics STM32MP157D-DK1 Discovery Kit}.

\subsection{Phase 1: Host System Preparation}
Build automation requires a Linux host environment (Ubuntu, Debian, Linux Mint, or WSL2/Docker).

\begin{cmdbox}{Step 1.1: Install Mandatory Host Packages}
First, update package indexes and install mandatory host compilation utilities:
\begin{lstlisting}[language=bash]
sudo apt update && sudo apt upgrade -y
sudo apt install -y build-essential libncurses5-dev libncursesw5-dev \
                    bison flex git rsync unzip bc wget cpio libssl-dev
\end{lstlisting}
\end{cmdbox}

\begin{cmdbox}{Step 1.2: Obtain Buildroot Source Code}
Create a dedicated working directory and clone the official Buildroot repository:
\begin{lstlisting}[language=bash]
mkdir -p ~/projects && cd ~/projects
git clone https://git.buildroot.net/buildroot
cd buildroot
\end{lstlisting}
\end{cmdbox}

\subsection{Phase 2: Target Configuration}

\begin{cmdbox}{Step 2.1: Inspect and Apply Board Defconfig}
Buildroot ships with pre-defined configurations inside the \texttt{configs/} directory. Inspect the available configurations for STM32 hardware:
\begin{lstlisting}[language=bash]
ls configs/ | grep stm32
\end{lstlisting}
Apply the default configuration for the STM32MP157 Discovery Kit:
\begin{lstlisting}[language=bash]
make st_stm32mp157a_dk1_defconfig
\end{lstlisting}
\textit{Note: This command generates a localized \texttt{.config} file containing all architecture, toolchain, kernel, and package parameters.}
\end{cmdbox}

\begin{cmdbox}{Step 2.2: Customize via Menuconfig}
Launch the ncurses graphical menu interface to review or modify build options:
\begin{lstlisting}[language=bash]
make menuconfig
\end{lstlisting}
Within \texttt{menuconfig}, engineers can configure:
\begin{itemize}[leftmargin=*]
    \item \textbf{Target Architecture:} ARM (little endian), Cortex-A7, NEON/VFPv4 floating point.
    \item \textbf{Toolchain:} Cross-compiler selection (GCC version, C library: \texttt{glibc}/\texttt{musl}/\texttt{uClibc}).
    \item \textbf{Kernel:} Linux kernel version, device tree sources (.dts).
    \item \textbf{Target Packages:} Adding lightweight tools (e.g., Python, dropbear SSH, i2c-tools).
    \item \textbf{Filesystem Images:} Ext4 rootfs, tarballs, and raw SD card image options.
\end{itemize}
Select \textbf{Exit} to return to the shell terminal.
\end{cmdbox}

\subsection{Phase 3: Cross-Compilation \& Build Artifact Inspection}

\begin{cmdbox}{Step 3.1: Execute Build}
Initiate the compilation process:
\begin{lstlisting}[language=bash]
make
\end{lstlisting}
\end{cmdbox}

\begin{tipbox}{Instructor Note: What Buildroot Performs During Compilation}
During \texttt{make}, Buildroot autonomously executes:
\begin{enumerate}[leftmargin=*]
    \item Downloads host build tools and target source archives.
    \item Cross-compiles the host cross-toolchain (\texttt{arm-linux-gnueabihf-gcc}).
    \item Compiles First Stage Bootloader (TF-A) and Second Stage Bootloader (U-Boot).
    \item Compiles the Linux Kernel image (\texttt{zImage}) and Device Tree Blob (\texttt{stm32mp157a-dk1.dtb}).
    \item Builds BusyBox and populates the root filesystem structure (\texttt{/bin}, \texttt{/etc}, \texttt{/lib}).
    \item Packs everything into output partition files and a raw \texttt{sdcard.img}.
\end{enumerate}
\end{tipbox}

\begin{cmdbox}{Step 3.2: Inspect Build Artifacts}
Once compilation completes, navigate to the output directory:
\begin{lstlisting}[language=bash]
ls -lh output/images/
\end{lstlisting}
Key files present:
\begin{itemize}[leftmargin=*]
    \item \texttt{sdcard.img}: Complete partitioned disk image ready to flash.
    \item \texttt{zImage}: Compressed Linux kernel binary.
    \item \texttt{stm32mp157a-dk1.dtb}: Compiled hardware Device Tree Blob.
    \item \texttt{rootfs.ext4}: Root filesystem partition file.
\end{itemize}
\end{cmdbox}

\subsection{Phase 4: Flashing Image to SD Card}

\begin{cmdbox}{Step 4.1: Identify Raw Storage Block Device}
Insert an SD card into the host computer. Identify its node using \texttt{lsblk}:
\begin{lstlisting}[language=bash]
lsblk
\end{lstlisting}
Identify the \textbf{raw block device name} (e.g., \texttt{/dev/mmcblk2} or \texttt{/dev/sdb}). Unmount any auto-mounted partition nodes:
\begin{lstlisting}[language=bash]
sudo umount /dev/mmcblk2p* 2>/dev/null || true
\end{lstlisting}
\end{cmdbox}

\begin{warningbox}{DANGER: Raw Disk Overwrite Risk with \texttt{dd}}
The \texttt{dd} utility performs bit-for-bit raw block copies. Specifying the wrong target drive (e.g., \texttt{of=/dev/sda}) will overwrite your host operating system drive, destroying all data without warning!
\begin{itemize}[leftmargin=*]
    \item Always target the raw disk node (\texttt{/dev/mmcblk2}), NOT a specific partition (\texttt{/dev/mmcblk2p1}).
    \item Double-check device size against \texttt{lsblk} output before pressing Enter.
\end{itemize}
\end{warningbox}

\begin{cmdbox}{Step 4.2: Flash Image with \texttt{dd}}
Execute raw bit-for-bit flashing to the target SD card:
\begin{lstlisting}[language=bash]
sudo dd if=output/images/sdcard.img of=/dev/mmcblk2 bs=1M status=progress
sync
\end{lstlisting}
Safely remove the SD card from the host.
\end{cmdbox}

\subsection{Phase 5: Hardware Target Connection \& Serial Console Setup}

\begin{conceptbox}{STM32MP157D-DK1 Hardware Configuration}
\begin{itemize}[leftmargin=*]
    \item \textbf{Boot Switches:} Set DIP switches \texttt{BOOT0} and \texttt{BOOT2} to \textbf{ON} to enforce booting from the MicroSD card slot.
    \item \textbf{Serial Console Connector:} Connect a Micro-USB cable from the host PC to the onboard ST-LINK USB port. The onboard ST-Link acts as a USB-to-UART bridge connected to UART4.
    \item \textbf{Power Supply:} Connect a dedicated 5V / 3A USB-C cable to the power connector port.
\end{itemize}
\end{conceptbox}

\begin{cmdbox}{Step 5.1: Configure Host User Serial Permissions}
Serial devices (\texttt{/dev/ttyACM0} or \texttt{/dev/ttyUSB0}) on Linux host machines are restricted to the \texttt{dialout} group. Add your host user to this group to prevent permission errors:
\begin{lstlisting}[language=bash]
sudo usermod -aG dialout $USER
groups $USER
\end{lstlisting}
\textit{Note: You must log out and log back in (or run \texttt{newgrp dialout}) for group changes to take effect.}
\end{cmdbox}

\begin{cmdbox}{Step 5.2: Open Serial Terminal with \texttt{picocom}}
Launch \texttt{picocom} at 115200 baud rate:
\begin{lstlisting}[language=bash]
picocom -b 115200 /dev/ttyACM0
\end{lstlisting}
Press the target board's \textbf{RESET} button. The serial terminal will display the full boot sequence: TF-A loader $\rightarrow$ U-Boot initialization $\rightarrow$ Linux Kernel boot log $\rightarrow$ init system startup.
\end{cmdbox}

\subsection{Phase 6: Target Bootup, Inspection, and Target Scripting}

Once booted, the Buildroot minimal OS presents a login prompt:
\begin{lstlisting}[language=bash]
buildroot login: root
\end{lstlisting}
\textit{Note: Default Buildroot images use username \texttt{root} with no password required.}

\begin{cmdbox}{Step 6.1: Explore Target Filesystem}
Inspect the system environment directly on the target hardware:
\begin{lstlisting}[language=bash]
# Check running kernel version
uname -a

# Inspect memory footprint
free -h

# Inspect mounted filesystems
df -h

# Verify CPU information
cat /proc/cpuinfo
\end{lstlisting}
Notice that the entire operating system consumes less than 30-50 MB of storage space and boots in under 3-5 seconds!
\end{cmdbox}

\begin{cmdbox}{Step 6.2: Create and Execute a Test Script on Target}
Embedded Linux provides standard POSIX shell scripting capabilities via BusyBox:
\begin{lstlisting}[language=bash]
vi test.sh
\end{lstlisting}
Press \texttt{i} to enter insert mode in \texttt{vi}, then add:
\begin{lstlisting}[language=bash]
#!/bin/sh
echo "---------------------------------------"
echo "  Hello from Custom Embedded Linux!   "
echo "  Board: STM32MP157D Discovery Kit    "
echo "---------------------------------------"
\end{lstlisting}
Press \texttt{Esc}, type \texttt{:wq} and press Enter to save and exit. Grant execution permissions and run the script:
\begin{lstlisting}[language=bash]
chmod +x test.sh
./test.sh
\end{lstlisting}
\end{cmdbox}

\begin{cmdbox}{Step 6.3: Exit Host Serial Terminal Session}
To disconnect from the target serial console in \texttt{picocom}, press:
\begin{center}
\textbf{\texttt{Ctrl + A}} followed by \textbf{\texttt{Ctrl + X}}
\end{center}
\end{cmdbox}
